%        File: Expose.tex
%     Created: Di Mrz 04 02:00  2014 Mitteleuropäische Z
% Last Change: Di Mrz 04 02:00  2014 Mitteleuropäische Z
%
% !TEX TS-program = pdflatex
% !TEX encoding = UTF-8 Unicode

\documentclass[11pt, titlepage=true]{scrartcl} % use larger type; default would be 10pt

\usepackage[utf8]{inputenc} % set input encoding (not needed with XeLaTeX)
\usepackage[ngerman]{babel} 
\usepackage[babel,german=quotes]{csquotes}
\usepackage[T1]{fontenc} % for enabling non-standard characters in
                         % hyphenation-list and smallcaps-boldfont

\usepackage[bibstyle=authortitle, citestyle=authoryear]{biblatex}
\addbibresource{Exposee}

% \DeclareUnicodeCharacter{00A0}{~} % for avoiding problems with no-break spaces

%%% PAGE DIMENSIONS
\usepackage{geometry} % to change the page dimensions
\geometry{a4paper} % or letterpaper (US) or a5paper or....
% \geometry{margin=2in} % for example, change the margins to 2 inches all round
% \geometry{landscape} % set up the page for landscape
%   read geometry.pdf for detailed page layout information

\usepackage{graphicx} % support the \includegraphics command and options

\graphicspath{{./Bilder/}}

\usepackage[parfill]{parskip} % Activate to begin paragraphs with an empty line rather than an indent

%%% PACKAGES
\usepackage{booktabs} % for much better looking tables
\usepackage{array} % for better arrays (eg matrices) in maths
%\usepackage{paralist} % very flexible & customisable lists (eg. enumerate/itemize, etc.)
\usepackage{verbatim} % adds environment for commenting out blocks of text & for better verbatim
\usepackage{subfig} % make it possible to include more than one captioned figure/table in a single float
% These packages are all incorporated in the memoir class to one degree or another...

%%% HEADERS & FOOTERS
\usepackage{fancyhdr} % This should be set AFTER setting up the page geometry
\pagestyle{fancy} % options: empty , plain , fancy
\renewcommand{\headrulewidth}{0pt} % customise the layout...
\lhead{}\chead{}\rhead{}
\lfoot{}\cfoot{\thepage}\rfoot{}

%%% SECTION TITLE APPEARANCE
\usepackage{sectsty}
\allsectionsfont{\sffamily\mdseries\upshape} % (See the fntguide.pdf for font help)
% (This matches ConTeXt defaults)

%%% ToC (table of contents) APPEARANCE
\usepackage[nottoc,notlof,notlot]{tocbibind} % Put the bibliography in the ToC
\usepackage[titles,subfigure]{tocloft} % Alter the style of the Table of Contents
\renewcommand{\cftsecfont}{\rmfamily\mdseries\upshape}
\renewcommand{\cftsecpagefont}{\rmfamily\mdseries\upshape} % No bold!

\usepackage[hidelinks]{hyperref}
\usepackage{ragged2e} %für links und rechtsbündige Umgebungen in Tabellen
\usepackage{todonotes}
\usepackage[normalem]{ulem} % ermöglicht es, Text mittels \sout und \xout
                            % durchzustreichen

%%% Eigene Befehle %%%
\newcommand{\was}[1]{\small\textit{#1}}

%%% Allgemeine Daten %%%
\newcommand{\autorOne}{Alexander Müller}
\newcommand{\matrikelnrOne}{940597}
\newcommand{\adresseOne}{Gartlager Weg 37, \\49086 Osnabrück}
\newcommand{\telOne}{01578 1901578}
\newcommand{\emailOne}{alemuell@uos.de}

\newcommand{\autorTwo}{János Sebestyén}
\newcommand{\matrikelnrTwo}{939525}
\newcommand{\adresseTwo}{Gartlager Weg 37, \\49086 Osnabrück}
\newcommand{\telTwo}{}
\newcommand{\emailTwo}{jsebesty@uos.de}

%%% Titlepage Expose %%%
\newcommand{\titel}{Ökologische Nachhaltigkeit durch \enquote{Nutzen statt Besitzen}?}
\newcommand{\stitle}{}
\newcommand{\art}{Exposé}
\newcommand{\fachgebiet}{Angewandte Systemwissenschaft}
\newcommand{\betreuung}{Katrin Bienge und Prof. Dr. Claudia Pahl-Wostl}
\newcommand{\institut}{Wuppertal Institut für Klima, Umwelt, Energie und Institut für Umweltsystemforschung}
\newcommand{\ort}{Osnabrück}

%%% END Article customizations

\begin{document}
\begin{titlepage}
  \begin{center}
    \includegraphics[scale=0.5]{Logo-Uni-Osnabrueck.jpg}\\[3ex]

    \vfill
    \LARGE{\textbf{\art}}\\[1.5ex]
    \huge{\textbf{\textsc{\titel}}}\\[1.0ex]
    \LARGE{\textbf{\stitle}}\\[6ex]
    \vfill

    \normalsize
    vorgelegt von:\\[12pt]

      \begin{minipage}[]{0.30\textwidth}
        \textbf{\autorOne}\\
        Matrikelnummer: \matrikelnrOne\\
        \emailOne\\
        \adresseOne\\
      \end{minipage}
      \hspace{0.15\textwidth}
      \begin{minipage}[]{0.3\textwidth}
        \textbf{\autorTwo}\\
        Matrikelnummer: \matrikelnrTwo\\
        \emailTwo\\
        \adresseTwo\\
      \end{minipage}

      \vfill
\todo{schön machen}
      \begin{minipage}[t]{0.45\textwidth}
        \begin{flushright}
          Datum:              \\[1.2ex]
          Ort:                \\[1.2ex]
          Betreuende Person:  \\[1.2ex]
          Institut:           \\[3ex]
        \end{flushright}
      \end{minipage}
      \hspace{0.01\textwidth}
      \begin{minipage}[t]{0.45\textwidth}
        \begin{flushleft}
          \today\\[1.2ex]
          \ort\\[1.2ex]
          \betreuung\\[1.2ex]
          \institut \\[3ex]
        \end{flushleft}
      \end{minipage}
  \end{center}
\end{titlepage}
\tableofcontents
\section{Thema der Arbeit}
Die heutigen industrialisierten Kulturen sind durch einen umfangreichen
individuellen Bedarf an Produkten und Dienstleistungen gekennzeichnet. Jeder
Haushalt, jede Person verfügt über zahlreiche, in der Summe höchst
ressourcenintensive Güter, was gesamtgesellschaftlich zu einer hohen, nicht mehr
tragfähigen Umweltbelastung führt.

Ein Problem heutiger Konsumkulturen, an dem Nachhaltigkeitsstrategien ansetzen
können, ist die ineffiziente Nutzung der produzierten Gütermenge: So wurde eine
durchschnittliche Bohrmaschine bei ihrer Entsorgung lediglich 50 Stunden genutzt, obwohl sie 250
Stunden hätte genutzt werden können\todo{Quelle und genauen Zahlen}. Autos sind im Durchschnitt mit nur 1,1 Personen besetzt, obwohl sie in der Regel Platz für vier Menschen bieten\todo{Quelle und
genauen Zahlen}. Diese Beispiele verdeutlichen, dass bei der individuellen Nutzung die maximale Nutzungsdauer und Auslastung von Produkten häufig nicht erreicht werden. Formen gemeinschaftlicher Nutzung, um die es in unserer Arbeit geht, können durch Nutzungsdauerverlängerung und Nutzungsintensivierung die daraus resultierenden Effizienzpotentiale heben und dadurch zu einer Reduktion der Umweltbelastung führen.

Zahlreiche dieser Nutzungsformen, die sich über alle Produktgruppen erstrecken,
sind bereits seit geraumer Zeit etabliert. Zum einen finden sich hier alle Formen des Gebrauchtwarenmarkts von Second-Hand-Läden für Kleidung über Geschäfte des An- und Verkaufs bis hin zu Flohmärkten. Zum anderen existiert eine breite Palette an Miet- und Verleihsystemen, wie zum
Beispiel Bibliotheken, Videotheken oder Autovermietungen. Neben diesen kommerziellen Angeboten gibt es auch vielfältige Formen der privaten gemeinschaftlichen Nutzung, wie
beispielsweise die gemeinsame Nutzung von Waschmaschinen, Rasenmähern oder Kühlschränken in
Wohnanlagen und Wohngemeinschaften, oder die Bildung von Fahrgemeinschaften.

Mit dem Aufkommen der heutigen Informations- und Kommunikationstechnologien
erhalten viele dieser Formen eine neue, ungeahnte Dimension: Das Internet
vergrößert die räumlichen Entfernungen und die Anzahl beteiligter
Teilnehmer*innen der herkömmlichen Gebraucht-, Miet- und Verleihmärkte um ein
Vielfaches. Durch das große Angebot und die einfache Möglichkeit der gezielten Suche verringert sich der Aufwand zum Zustandekommen des Warenaustauschs. Ermöglicht durch den einfachen, ortsungebundenen Zugang zum
Internet und die Vernetzung von Menschen mit ähnlichen Bedürfnissen entstanden ganz neue
Formen der gemeinschaftlichen Nutzung: So ermöglichen es entsprechende
Internetplattformen beispielsweise, anstatt in einem Hotel zu übernachten,
auf einer Reise ganz bequem bei Privatpersonen unterzukommen oder
mit Hilfe eines Smartphones ganz einfach das nächste Auto eines Floating
Car Sharing Systems \todo{richtigen Begriff rausfinden} zu finden.

Hinter diesen neuen Formen werden unter Verwendung der Begriffe
\enquote{collaborative consumption} oder \enquote{Kokonsum} soziale Innovationen
gesehen, die zu mehr Nachhaltigkeit bei gleichbleibendem oder sogar gesteigertem
Wohlstand führen können. Und aus Unternehmensperspektive stellen sich diese neuen
Nutzungsformen als wichtiger neuer
ökonomischer Wachstumsmarkt -- die \enquote{Sharing Economy} -- dar. In der Forschung wird das Thema unter den Begriffen \enquote{Nutzen statt Besitzen}, \enquote{eigentumsersetzender Konsum}, \enquote{eco services} und \enquote{product-service systems} diskutiert. \todo{Begriffe definieren}

Im Vergleich zur Euphorie, die von der Sharing Economy ausgeht, zeigt sich aus wissenschaftlicher Perspektive hinsichtlich der Umweltwirkungen gemeinschaftlicher Nutzungssysteme ein differenzierteres Bild: Zum einen werden die positiven Primäreffekte der
Nutzungsdauerverlängerung und -intensivierung um weitere sekundäre
Ressourceneinsparmöglichkeiten ergänzt. So kann es beispielsweise bei der
gemeinschaftlichen Nutzung einer Waschmaschine sinnvoll sein, eine
halbgewerbliche Maschine mit entsprechend längerer Lebensdauer und höherer
Effizienz zu verwenden. Zum anderen werden die durch gemeinschaftliche Nutzung
ermöglichten Ressourceneinsparungen jedoch auch durch negative Umweltwirkungen
begleitet: Beispielsweise müssen für die Koordination der gemeinschaftlichen
Nutzung und den Transport der genutzten Produkte zwischen den Nutzer*innen zusätzliche
Ressourcen aufgewendet werden. Und auch positive sekundäre Effekte können mit negativen Nebenwirkungen einhergehen, wie beispielsweise für die Produktion einer halbgewerblichen Waschmaschine im Vergleich zu einer handelsüblichen
mehr Ressourcen benötigt werden.

Nicht zuletzt wirken sich die aus der Perspektive der Sharing Economy positiven
finanziellen Einsparungen bei der gemeinschaftlichen Nutzung auf die Nachfrageseite der
Konsumenten aus: Das eingesparte Geld kann in weiteren Konsum und damit
zusätzlichen Ressourcenverbrauch investiert werden (Rebound-Effekt). Außerdem
bleiben in einer solchermaßen veränderten Ökonomie, welche den Zugang und nicht
den Besitz in den Vordergrund stellt, die Bedürfnisstrukturen nicht unberührt:
Die Möglichkeit, auf viele (Luxus-)Güter zuzugreifen, könnte mit einem
verstärkten Wunsch einhergehen, dies auch tatsächlich zu tun.

Welche Netto-Umweltwirkungen schließlich durch die gemeinschaftliche Nutzung zu
erwarten sind, hängt in hohem Maße von dem betrachteten Produkt und dem konkreten
Nutzungssystem ab. Entscheidende Produktmerkmale sind zum Beispiel, ob während der Nutzung des Produkts Ressourcen verbraucht werden (wie bei Elektrogeräten), oder wie groß das Verhältnis zwischen typischer individueller und maximal möglicher Nutzungsdauer ist (im oben genannten Beispiel der Bohrmaschine 1:5). Zu den wichtigen Eigenschaften des Nutzungssystems zählen eventuell anfallende zusätzliche Transporte (wie beim Second-Hand-Handel) und die Auslastung während der Nutzung (zum Beispiel 3 Personen pro Auto). Die Zusammenhänge zwischen Merkmalen der Produkte und Nutzungssysteme und den resultierenden Umweltwirkungen zu analysieren ist der Gegenstand unserer Arbeit.

\section{Forschungsstand und Forschungslücke}\label{sec:forschung}
\was{An welchen derzeitigen Forschungsstand knüpfen wir mit unserem Thema an? Hier führen wir relevante Forschung und aktuelle Studien an.}

- Beschreibung / Einordnung der Literatur, die wir schon gelesen haben

Zu den Umwelteffekten kollektiver Nutzungsformen wurde bereits einiges
publiziert: 
\begin{enumerate}
     \item IÖW-Studie
       \begin{itemize}
         \item In den Jahren 1997-2000 bearbeitete das Institut für ökologische
           Wirtschaftsforschung (IÖW) das vom Bund geförderte
           Forschungsvorhaben \enquote{Ökologische Entlastungspotentiale,
             Umsetzungsprobleme und Entwicklungsperspektiven von Strategien zur
             Nutzungsdauerverlängerung (NV) und Nutzungsintensivierung (NI).}
             \textcite{scholl_produkte_1998} stecken zunächst den
             begrifflichen Rahmen für die Beschreibung gemeinschaftlicher
             Nutzungsformen und deren ökologischen Bewertung ab. Hieran
             schließen sich eine Fallstudie zur privaten Textilwäsche
             \parencite{hirschl_produkte_2000} und eine zum Wintersport
             \parencite{konrad_produkte_2000} an, wobei hier insbesondere der
             Frage nachgegangen wird, wie das ökologische Entlastungspotenzial
             von Nutzungskonzepten im Vergleich zur herkömmlichen individuellen
             Nutzung einzuschätzen ist. Eine zusammenfassende Darstellung der
             Projektergebnisse geben schließlich
             \textcite{hirschl_nachhaltige_2001}. Aus einem Vergleich der beiden
             Fallstudien leiten sie allgemeinere Erkenntnisse für die
             Nachhaltigkeitsstrategien Nutzungsdauerverlängerung und
             -intensivierung ab.
       \end{itemize}

     \item EU-Studie
       \begin{itemize}
         \item In einem ähnlichen Zeitrahmen wie die Studie von dem IÖW wurde
           auch von verschiedenen weiteren
           Forschungseinrichtung ein EU finanziertes Forschungsprojekt
           unter dem Titel \enquote{Eco-Services for Sustainable Development in
           the European Union} durchgeführt, welche einen sehr ähnlichen Aufbau
           besitzt, allerdings zum Teil andere Fallstudien betrachtet und durch
           den internationalen Kontext auch den Blick auf länderspezifische
           Unterschiede wirft. 
           Die vertiefenden Fallstudien, welche ökologische Betrachtungen
           einschließen, wurden in den  Bereichen Heimwerken, Baueigenleistung
           und Gartenpflege \cite{behrendt_oko-rent_2000}*, nachhaltiges
           Waschen \cite{behrendt_nachhaltig_2000}* und
           Car-Sharing \cite{behrendt_car-sharing_2000}* durchgeführt. Eine
           Systematisierung der betrachteten Produktdienstleistung und eine
           vergleichende Darstellung der Fallstudien, Zusammenfassung des
           Projekts geben \textcite{behrendt_eco-service_2003}.
       \end{itemize}
     \item weitere Studien
       \begin{itemize}
         \item In der jüngeren Literatur finden sich weitere Publikationen,
           welche sich zum einen theoretisch dem Thema kollektiver
           Nutzungsformen zuwenden
           \parencites{scholl_marketing_2009}{scholl_nutzen_2012} und zum
           anderen empirisch-analytisch die Umwelteffekte ausgewählter
           Produkt-Nutzungs-Systeme
           erheben
           \parencites{rabelt_nachhaltiger_2007}{erdmann_quantifizierung_2011}.
           Erwähnenswert ist, dass \textcite{scholl_marketing_2009} eine
           umfassende Liste der bei kollektiven Nutzungsformen 
           auftretenden positiven und negativen Umwelteffekte aufstellt. Er
           gliedert diese ein in solche, welche sich direkt durch die
           Nutzungsform ergeben und solche, welche indirekt durch veränderte
           Nachfragemuster resultieren.
           Während viele der bisher zitierten Studien letztere Gruppe von
           Effekten nicht
           näher quantifizieren, versucht
           \textcite{erdmann_quantifizierung_2011}  genau das anhand des
           online-basierten Gebrauchtwarenhandels.
       \end{itemize}
     \item ökonomische Modelle zum Thema Wiederverwendung
       \begin{itemize}
         \item Auch in der ökonomischen Literatur lassen sich Autoren und
           Autorinnen finden, die sich mit kollektiven Nutzungsformen
           unter Verwendung einer deutlich formalisierteren Methodik
           auseinandersetzen. \textcite{yokoo_economic_2010} entwickelt ein
           ökonomisches Modell der Wiederverwendung um Wohlfahrts- und
           Abfalleffekte auf Second-Hand-Märkten zu untersuchen.
           \textcite{thomas_environmental_2011} entwickelt ein ähnliches
           ökonomisches Modell, fokussiert jedoch, in dem er dieses in ein Modell für
           eine Lebenszyklusbilanz integriert, noch stärkere den Blick auf die
           ökologischen Effekte auf Second-Hand-Märkten. Zur Illustration des
           Modells wird hier der Gebrauchtwarenhandel mit Büchern betrachtet.       \end{itemize}
\end{enumerate}

Die angegebenen Kurzzusammenfassungen der ausgewählten Literatur beziehen sich
vornehmlich auf die für unsere Arbeit relevanten Inhalte. Darüber hinaus sind
die zitierten Quellen meistens viel umfassender angelegt. So werden --
neben den Umweltwirkungen --  häufig noch
Handlungsempfehlungen, Strategien zur Förderung, Hemmnisse,
Marktanalysen, Anschlussfähigkeit, Typisierung von Konsument*innen und die
Motivation für solche Nutzungsformen betrachtet.

- Zusammenfassung / Fazit / Forschungslücke (s.u.)

Insgesamt kann festgestellt werden, dass zu dem Thema kollektive Nutzungsformen
bereits einiges sowohl an empirischer als auch konzeptioneller Forschung
betrieben wurde. Die Quantifizierung von Umwelteffekten erfolgt in den zitierten
Quellen immer an speziellen Produkten (Werkzeuge, textiles
Waschen, ...) und Nutzungssystemen (Gebrauchtwarenhandel übers Internet,
Car-Sharing). 
 
\was{Gibt es eine Forschungslücke im derzeitigen Forschungsstand? Wie werden wir
 diese mit unserer Arbeit schließen?}
 - immer nur detaillierte Betrachtung einzelner Produkte
 - kein einheitliches methodisches Vorgehen: Bezugsgröße variiert
 - Einzelfallberechnungen, keine Analyse der Parameter, Potentiale
 - nicht allgemein
 - keine konsequente Betrachtung der Ressourceninputs
 
 
\section{Forschungsziel und Fragestellung}
Das Hauptziel unserer Arbeit ist die Entwicklung von Kriterien für die Senkung des Umweltverbrauchs durch gemeinschaftliche Produktnutzung. Darunter fällt eine systematische und möglichst allgemeine Betrachtung der verschiedenen Umweltwirkungen von gemeinschaftlichen Nutzungsformen. Ein methodologisches Ziel der Arbeit ist es, einen einheitlichen Methodenrahmen für die Untersuchung verschiedener Produkte und Nutzungssysteme zu entwickeln.

Eine mit diesem Forschungsziel korrespondierende Fragestellung lautet:
\begin{itemize}
	\item Unter welchen Umständen kann der Umweltverbrauch eines Produktes durch gemeinschaftliche Nutzung gegenüber der individuellen Nutzung gesenkt werden?
\end{itemize}
Etwas feinteiliger betrachtet, lässt sich diese Fragestellung in die folgenden zwei Teilfragen unterteilen:
\begin{itemize}
	\item Was sind die Mechanismen, die den Umweltverbrauch bei gemeinschaftlicher Nutzung bestimmen und wie wirken diese?
	\item Welche Eigenschaften des Produkts und der Nutzungsform beeinflussen die Wirkung dieser Mechanismen und wie fließen sie ein?
\end{itemize}

\section{Methodische Vorgehensweise}
Eine Grundannahme, unter der wir an die Bearbeitung der Fragestellung herangehen, ist, dass der Umweltverbrauch eines Produktnutzungssystems durch Merkmale des Produktes (einschließlich seines Produktions- und Entsorgungssystems) und Merkmale des Nutzungssystems bestimmt wird. Ziel der Untersuchung ist es also, den Zusammenhang zwischen diesen Merkmalen und dem Umweltverbrauch zu bestimmen.

Die Zielgröße \enquote{Umweltverbrauch} operationalisieren wir in unserer Untersuchung durch die Menge der eingesetzten Ressourcen und der ausgestoßenen Emissionen. Für die Ressourcenseite werden wir auf das MIPS-Konzept (\cite{schmidt-bleek_maia:_1998}) zurückgreifen. Bei der Betrachtung der Emissionen beschränken wir uns auf Treibhausgase (CO$_{2}$-Äquivalente). Je nach Wirkungsmechanismus wird es erforderlich sein, unterschiedliche Betrachtungsrahmen für die Bilanzierung heranzuziehen. So können die Umwelteffekte beispielsweise auf eine einzelne Produktnutzung, auf den gesamten Konsum einer Person oder gar auf den Konsum einer ganzen Volkswirtschaft bezogen werden. Wir streben eine möglichst einheitliche Betrachtung an, die allerdings an ihre Grenzen stoßen wird, da beispielsweise Effekte einer gemeinschaftlichen Produktnutzung auf andere Konsumbereiche mit einer Produkt-zentrierten Betrachtung nicht erfasst werden können.

Um die mögliche Senkung des Umweltverbrauchs durch gemeinschaftliche Nutzung zu berechnen, wird als Referenzszenario jeweils die individuelle Nutzung herangezogen. Ergänzend dazu kann das Resultat auch mit den Effekten durch andere Maßnahmen zur Senkung des Umweltverbrauchs ins Verhältnis gesetzt werden, sofern dies sinnvoll erscheint. Beispielsweise könnten technische Effizienzsteigerungen oder Konsumreduzierungen betrachtet werden.

Die in Tabelle \ref{tab:Mechanismen} aufgeführten Wirkungsmechanismen möchten wir prinzipiell in unsere Betrachtung miteinbeziehen. Die Liste basiert auf \cite{scholl_marketing_2009}, wobei wir einige der dort aufgeführten Mechanismen aus Plausibilitätsüberlegungen nicht aufgenommen und andere selbst hinzugefügt haben.

{\tiny \begin{table}[h]
	\begin{tabular}{p{7cm}p{7cm}}
		 \toprule
		 \multicolumn{2}{l}{\textbf{Umweltauswirkungen durch die Nutzung}} \vspace{0.2 cm} \\
		 \textbf{positiv} & \textbf{negativ} \\
		 \midrule
		 Nutzungsdauer\-verlängerung \vspace{0.2cm} &  Zusätzliche Transaktionen \vspace{0.2cm} \\
		 Einsatz langlebiger Produkte \vspace{0.2cm} &  zusätzlicher Ressourcenverbrauch für Langlebigkeit \vspace{0.2cm} \\
		 Verwendung verbrauchsarmer bzw. leistungsstarker Geräte \vspace{0.2cm} &  Beschleunigte Ausmusterung \vspace{0.2cm} \\
		 Maximierung der Geräteauslastung \vspace{0.2cm} & Tauschbedingter Verschleiß \vspace{0.2cm} \\
		 Berücksichtigung des technischen Fortschritts \vspace{0.2cm} & \vspace{0.2cm} \\
		 Wartung / Reparaturen \vspace{0.5cm}& \\
		 \midrule
		 \multicolumn{2}{l}{\textbf{Umweltauswirkungen durch Nachfrageänderung}} \vspace{0.2 cm} \\
		 \textbf{positiv} & \textbf{negativ} \\
		 \midrule
		 Nachfrageverringerung durch höhere Kostentransparenz \vspace{0.2cm} & Erleichterter Produktzugang \vspace{0.2cm} \\
		 Vermeidung von Fehlkäufen \vspace{0.2cm} & Wunsch nach Eigentum \vspace{0.5cm} \\
		  & Rebound-Effekt \vspace{0.2cm} \\
		 \bottomrule
	\end{tabular}
	\caption{Betrachtete Mechanismen für die Umweltauswirkungen durch gemeinschaftliche Nutzung. Quelle: verändert nach \cite{scholl_marketing_2009}.}
	\label{tab:Mechanismen}
\end{table}}

Nicht alle Mechanismen sind für alle Produkte und Nutzungssysteme relevant. Welche Mechanismen für eine bestimmte Produkt-Nutzungssystem-Kombination zum Tragen kommen, hängt von spezifischen Annahmen über die Produkte und Nutzungssysteme ab, die im Laufe der Arbeit getroffen werden. Auch sei erwähnt, dass nicht alle Mechanismen ausschließlich bei gemeinschaftlichen Nutzungsformen wirken, sondern auch bei individueller Nutzung auftreten können.

Um die Mechanismen und deren Auswirkungen quantifizierbar und vergleichbar zu machen, werden wir ein Modell entwickeln. Der Modellzweck ist, die Zusammenhänge zwischen den Merkmalen der Produkte und Nutzungssysteme und dem Umweltverbrauch abzubilden. Während das Modell letztlich möglichst allgemein anwendbar sein soll, lassen wir uns bei der Modellentwicklung durch zwei Fallstudien leiten. Sie dienen der Veranschaulichung der abstrakten Zusammenhänge und gewährleisten die praktische Anwendbarkeit des Modells. Zudem kann die Fokussierung auf konkrete Fallstudien die Modellanalyse vereinfachen, da der potentiell sehr große Parameterraum etwas eingeschränkt werden kann. Auswahlkriterien für die Fallstudien waren die Verfügbarkeit von Literatur und die Abdeckung möglichst unterschiedlicher Produkteigenschaften, damit alle potentiellen Mechanismen zur Anwendung kommen. Die Entscheidung fiel auf die zwei Produkte Waschmaschinen und Bücher.

Konkret wird unsere Arbeit durch den folgenden groben Ablauf geleitet sein:
\begin{itemize}
	\item Modellentwicklung für Einzelprozesse
	\item Analyse der Einzelmodelle
	\item Modellkopplung mehrerer wechselwirkender Prozesse
	\item Analyse der gekoppelten Modelle
	\item Anwendung der Einzelmodelle und gekoppelten Modelle auf Fallstudien
\end{itemize}


\newpage
\section{Aufbau der Arbeit}
Der Aufbau der Arbeit wird sich an folgender Gliederung orientieren:

\begin{enumerate}
\item Einleitung
	\begin{enumerate}
		\item Thema
		\item Forschungsstand
		\item Forschungsziel und Fragestellung
	\end{enumerate}
\item Methodisches Vorgehen
	\begin{enumerate}
		\item Modellansatz
		\item Fallstudien
	\end{enumerate}
\item Modell
	\begin{enumerate}
		\item Einzeleffekte
		\item Effekt-Kopplungen
		\item Modellübersicht
	\end{enumerate}

\item Synthese
\item Ausblick/Diskussion
\item Literaturverzeichnis
\item Anhang
\end{enumerate}

\section{Zeitplan und einzelne Arbeitsschritte}
\newcommand{\paketEins}{Literaturarbeit}
\newcommand{\paketZwei}{Einzelmodelle erstellen und analysieren}
\newcommand{\paketDrei}{Modellkopplungen erstellen und analysieren}
\newcommand{\paketVier}{Text verfassen}
Die Bearbeitung des vorgestellten Themas lässt sich grob in die vier
Arbeitspakete \paketEins, \paketZwei, \paketDrei\ und \paketVier\ einteilen. 

Das Arbeitspaket \enquote{\paketEins}\ umfasst die abschließende Sichtung und
Zusammenfassung der im Abschnitt \ref{sec:forschung} vorgestellten Literatur. Da
die Literatur für die Erstellung des Exposees bereits ausgewählt und zum Teil
auch bearbeitet wurde, ist eine Fertigstellung dieses Arbeitspakets bis Ende des
Jahres geplant.

Das von der Planung her umfangreichste Paket \enquote{\paketZwei}\ ist der Kern
unserer Master-Arbeit. Der geschätzte Umfang ergibt sich daraus, dass relativ
viele, zum Teil sehr unterschiedliche Mechanismen für die Umweltauswirkungen
durch gemeinschaftliche Nutzung betrachtet werden und die Formulierung,
Implementierung und Analyseform der Modelle noch ungewiss sind.

Vorausgesetzt dass sich diese Unsicherheiten mit Abschluss des zweiten
Arbeitspakets klären, erwarten wir, dass der Umfang des sich anschließenden
Arbeitspakets \enquote{\paketDrei}\ geringer ist. Hier wird die Schwierigkeit
vor allem in
einer sinnvollen Auswahl von gekoppelten Modellen bestehen. Eventuell müssen die
Analysen der Einzemodelle für die Modellkopplungen angepasst werden, die
grundsätzliche Struktur kann unserer Erwartung nach jedoch übernommen werden.

Im Arbeitspaket \enquote{\paketVier}\ wird schließlich das Augenmerk auf die
Erstellung eines zusammenhängenden Dokuments gelenkt. Die abschließende
schriftliche Dokumentation der vorangegangenen Arbeitspakte wird hier ebenso
Bestandteil sein, wie die Formulierung einer Einleitung und eines
Synthese-Kapitels, die Diskussion der Ergebnisse, sowie die sprachliche
Überarbeitung und die Formatierung des Textes.

Während der Vorlesungszeit planen wir im Umfang von ca. 20 Personenstunden pro
Woche an der Master-Arbeit zu arbeiten. In der vorlesungsfreien Zeit würde sich
dann der Arbeitsumfang auf ca. 70 Personenstunden pro Woche erhöhen. Zusammen
mit dem geschätzten Aufwand je Arbeitspaket ergibt sich der voraussichtliche
Abschluss der Masterarbeit bis zum 31.05.2015 (vgl. Tabelle \ref{tab:zeitplan}).
Wir planen eine maximale Verzögerung des Projekts im Umfang von 25\% des
geschätzten Gesamtaufwands ein. Der späteste Termin für die Beendigung des
Projekts ist demnach der 12.07.2015.

\begin{table}[h]
  \centering
  \begin{tabular}{lccc}
    \toprule
    Arbeitspaket    & Umfang\textsuperscript{*} & Beginn      & Ende      \\
    \midrule
    \paketEins      & 80                        & 24.11.14    & 17.12.14  \\
    \paketZwei      & 500                       & 05.01.15    & 22.03.15  \\
    \paketDrei      & 280                       & 23.03.15    & 19.04.15  \\
    \paketVier      & 420                       & 20.04.15    & \textbf{31.05.15}\\
    \midrule
    Puffer          & 420                       & 01.06.15    & 12.07.15  \\
    \bottomrule
  \end{tabular}\\
  \footnotesize \textsuperscript{*}Umfang angegeben in
  Personenstunden\hspace{7.5cm}
  \caption{Grobe Strukturierung der Master-Arbeit in vier Arbeitsschritte mit
    Aufwandsschätzung}
  \label{tab:zeitplan}
\end{table}

\section{Auswahlbibliographie}
\cite{tukker_eight_2004}
\printbibliography
\end{document}


